% Compilar desde este directorio con:
% pdflatex Clase_06_Ingenieria_de_caracteristicas.tex
\documentclass[aspectratio=169,10pt]{beamer}

\usepackage[utf8]{inputenc}
\usepackage[T1]{fontenc}
\usepackage[spanish,es-nodecimaldot]{babel}
\usepackage{amsmath,amssymb,booktabs,tabularx,array,tikz}
\usetikzlibrary{arrows.meta,positioning,calc,shapes.geometric}

\definecolor{UAPNavy}{HTML}{17324D}
\definecolor{UAPTeal}{HTML}{007F82}
\definecolor{UAPGold}{HTML}{E2A33A}
\definecolor{UAPLight}{HTML}{F3F6F8}
\definecolor{UAPGray}{HTML}{536471}
\definecolor{UAPRed}{HTML}{A94343}

\tikzset{
  flow/.style={-{Stealth[length=2.2mm]},thick,draw=UAPNavy},
  datum/.style={draw=UAPGold,fill=UAPGold!16,rounded corners=2pt,align=center,minimum height=8mm,text=UAPNavy},
  process/.style={draw=UAPTeal,fill=UAPTeal!12,rounded corners=3pt,very thick,align=center,minimum height=10mm,text=UAPNavy},
  decision/.style={draw=UAPNavy,fill=UAPNavy!9,rounded corners=3pt,very thick,align=center,minimum height=10mm,text=UAPNavy}
}

\usetheme{Boadilla}
\usecolortheme{default}
\setbeamercolor{normal text}{fg=UAPNavy,bg=white}
\setbeamercolor{structure}{fg=UAPTeal}
\setbeamercolor{frametitle}{fg=white,bg=UAPNavy}
\setbeamercolor{title}{fg=white}
\setbeamercolor{subtitle}{fg=UAPGold}
\setbeamercolor{block title}{fg=white,bg=UAPTeal}
\setbeamercolor{block body}{fg=UAPNavy,bg=UAPLight}
\setbeamercolor{alertblock title}{fg=white,bg=UAPRed}
\setbeamercolor{alertblock body}{fg=UAPNavy,bg=UAPRed!7}
\setbeamerfont{frametitle}{series=\bfseries}
\setbeamertemplate{navigation symbols}{}
\setbeamertemplate{itemize item}{\color{UAPGold}$\blacktriangleright$}
\setbeamertemplate{itemize subitem}{\color{UAPTeal}$\bullet$}
\setbeamertemplate{footline}{%
  \leavevmode\hbox{%
    \begin{beamercolorbox}[wd=.78\paperwidth,ht=2.8ex,dp=1.2ex,leftskip=1em]{author in head/foot}%
      \color{white}\insertshorttitle
    \end{beamercolorbox}%
    \begin{beamercolorbox}[wd=.22\paperwidth,ht=2.8ex,dp=1.2ex,rightskip=1em plus1fil]{date in head/foot}%
      \color{UAPGray}\hfill Semana 6 \enspace | \enspace \insertframenumber/\inserttotalframenumber
    \end{beamercolorbox}%
  }%
}

\newcommand{\concepto}[1]{\textcolor{UAPTeal}{\textbf{#1}}}
\newcommand{\br}{\\}

\title[IA · Clase 6]{Ingeniería de características}
\subtitle{De datos brutos a representaciones útiles}
\author{Curso de Inteligencia Artificial}
\institute{Semana 6 · Clase 6}
\date{}

\begin{document}

% 1
\begin{frame}[plain]
  \begin{beamercolorbox}[wd=\paperwidth,ht=\paperheight,sep=2em]{frametitle}
    \vfill
    {\usebeamerfont{title}\usebeamercolor[fg]{title}\Huge\inserttitle\par}
    \vspace{.45em}{\usebeamerfont{subtitle}\usebeamercolor[fg]{subtitle}\Large\insertsubtitle\par}
    \vspace{1em}
    \begin{tikzpicture}[node distance=5mm,scale=.76,transform shape]
      \node[datum] (d) {Dato bruto}; \node[process,right=of d] (c) {Característica};
      \node[process,right=of c] (p) {Pipeline}; \node[decision,right=of p] (m) {Modelo};
      \draw[flow] (d)--(c); \draw[flow] (c)--(p); \draw[flow] (p)--(m);
    \end{tikzpicture}\par
    \vspace{.7em}{\color{white}\small El modelo solo puede aprender diferencias que su representación vuelve visibles.\par}
    \vfill {\color{UAPGold}\large\insertauthor\par}{\color{white}\insertinstitute\par}
  \end{beamercolorbox}
\end{frame}

% 2
\begin{frame}{La representación define qué puede aprenderse}
  \footnotesize
  Una \concepto{característica} o \concepto{atributo} es una magnitud utilizada como entrada de un modelo.
  \begin{tabularx}{\textwidth}{p{.24\textwidth} p{.29\textwidth} X}
    \toprule Dato disponible & Representación posible & Supuesto incorporado \\ \midrule
    fecha y hora & hora, día, feriado & existe un patrón de calendario \\
    lluvia acumulada & valor, nivel o indicador & importa la magnitud o un umbral \\
    ubicación & zona, distancia o conectividad & la cercanía relevante está bien definida \\
    \bottomrule
  \end{tabularx}
  \begin{block}{Criterios}
    Una característica útil tiene significado, está disponible al predecir, se calcula igual durante entrenamiento e inferencia y mantiene estabilidad operativa.
  \end{block}
\end{frame}

% 3
\begin{frame}{Características numéricas}
  \footnotesize
  \begin{tabularx}{\textwidth}{>{\bfseries}p{.2\textwidth} p{.28\textwidth} X}
    \toprule Operación & Uso & Riesgo \\ \midrule
    Estandarizar & comparar escalas & aprender media y dispersión con todos los datos \\
    Logaritmo & representar cambios relativos & aplicarlo a ceros o negativos \\
    Razón & normalizar por exposición & denominadores pequeños o nulos \\
    Recorte & limitar extremos & ocultar eventos reales \\
    \bottomrule
  \end{tabularx}
  \vspace{.35em}
  Si una zona registra 120 viajes con capacidad 150:
  \[
    \text{presión}=120/150=0{,}80.
  \]
  \begin{alertblock}{Control}
    La razón facilita comparar zonas de tamaños distintos, pero necesita una regla para capacidad cero.
  \end{alertblock}
\end{frame}

% 4
\begin{frame}{Características categóricas}
  \small
  \begin{itemize}
    \item Una categoría nominal no tiene distancia ni orden intrínsecos.
    \item Codificar \texttt{norte}, \texttt{centro}, \texttt{sur} como 1, 2, 3 crea una geometría artificial.
    \item La codificación indicadora crea una columna binaria por nivel.
    \item Las categorías ordinales admiten orden, pero no necesariamente intervalos iguales.
    \item Los niveles raros y desconocidos necesitan una política explícita.
  \end{itemize}
\end{frame}

% 5
\begin{frame}{Categoría desconocida: ejemplo y cuidados}
  \footnotesize
  \begin{block}{Caso}
    El modelo se entrenó con \texttt{zona} $\in \{\texttt{norte},\texttt{centro},\texttt{sur}\}$ y durante el uso aparece \texttt{oeste}.
  \end{block}
  \begin{alertblock}{Tratamiento riesgoso}
    Un pipeline configurado para sustituir valores no reconocidos por la categoría más frecuente podría codificar \texttt{oeste} como \texttt{norte}. Esto evita un error técnico, pero no demuestra que ambas zonas se comporten igual y oculta una posible diferencia entre entrenamiento y uso.
  \end{alertblock}
  \textbf{Cuidados necesarios}
  \begin{itemize}
    \item distinguir una categoría nueva de un dato faltante;
    \item comprobar si proviene de un error, una fuente nueva o una ampliación de la población;
    \item representarla como \texttt{desconocida} u \texttt{otra} con un codificador preparado;
    \item registrar su frecuencia y revisar el desempeño del modelo en esos casos;
    \item actualizar y reevaluar el modelo si la categoría se vuelve habitual.
  \end{itemize}
\end{frame}

% 6
\begin{frame}{Tiempo y espacio}
  \small
  Para una posición temporal $t$ dentro de un periodo $P$:
  \[
    x_{\sin}=\sin(2\pi t/P),\qquad x_{\cos}=\cos(2\pi t/P).
  \]
  El par conserva que las 23:00 y las 00:00 son horas cercanas. Los rezagos y medias móviles deben usar únicamente valores anteriores al instante de predicción.
  \vspace{.55em}

  En espacio pueden utilizarse distancias, pertenencia a zonas, accesibilidad o conectividad. La distancia geométrica no siempre representa tiempo de viaje ni transferencia del fenómeno.
\end{frame}

% 7
\begin{frame}{Interacciones y conocimiento del dominio}
  \footnotesize
  Una interacción representa que el efecto de una variable depende de otra:
  \[
    f(x)=\beta_0+\beta_1x_1+\beta_2x_2+\beta_{12}x_1x_2.
  \]
  \textbf{Ejemplo:} lluvia $\times$ hora pico puede representar que la lluvia afecta más la circulación cuando la red está congestionada.
  \begin{itemize}
    \item Priorizar combinaciones con una explicación operativa.
    \item Verificar que existan suficientes observaciones de la combinación.
    \item Comparar su aporte fuera de muestra.
  \end{itemize}
  \begin{alertblock}{Escala del problema}
    Con 100 variables existen 4950 interacciones por pares: generar todas aumenta capacidad y riesgo de sobreajuste.
  \end{alertblock}
\end{frame}

% 8
\begin{frame}{Selección de atributos}
  \footnotesize
  Seleccionar atributos conserva un subconjunto de las variables disponibles.
  \begin{tabularx}{\textwidth}{>{\bfseries}p{.18\textwidth} p{.3\textwidth} X}
    \toprule Familia & Estrategia & Limitación principal \\ \midrule
    Filtro & puntuar antes del modelo final & puede ignorar interacciones \\
    Envolvente & comparar subconjuntos con el modelo & costo y sobreajuste de la búsqueda \\
    Embebida & seleccionar durante el ajuste & depende del modelo y la escala \\
    \bottomrule
  \end{tabularx}
  \vspace{.4em}
  Relevancia no equivale a causalidad. Dos variables correlacionadas pueden sustituirse y cambiar de posición entre muestras.
  \begin{block}{Regla}
    Cualquier selección basada en datos se ajusta dentro de cada fold.
  \end{block}
\end{frame}

% 9
\begin{frame}{Métodos de filtro}
  \footnotesize
  Los filtros puntúan atributos sin entrenar repetidamente el modelo final.
  \begin{tabularx}{\textwidth}{>{\bfseries}p{.2\textwidth}p{.31\textwidth}X}
    \toprule Criterio & Qué busca & Ejemplo \\ \midrule
    Varianza & columnas constantes o casi constantes & retirar un sensor que siempre informa 1 \\
    Asociación & correlación lineal, relación monótona o asociación categórica con $y$ & ordenar velocidad por su asociación con bloqueo \\
    Información mutua & dependencia lineal o no lineal con $y$ & detectar señal aportada por lluvia \\
    \bottomrule
  \end{tabularx}
  \begin{block}{Ejemplo ilustrativo}
    Información mutua con bloqueo: velocidad $=0{,}31$, lluvia $=0{,}18$, hora $=0{,}05$ e identificador $=0{,}01$. Un selector \textit{top-2} conservaría velocidad y lluvia.
  \end{block}
  \begin{alertblock}{Cuidados}
    El número de atributos se valida; un filtro univariado puede perder interacciones y todo filtro que usa $y$ se ajusta dentro de cada fold.
  \end{alertblock}
\end{frame}

% 10
\begin{frame}{Métodos envolventes}
  \footnotesize
  Evalúan subconjuntos entrenando el modelo que finalmente se utilizará.
  \begin{tabularx}{\textwidth}{>{\bfseries}p{.27\textwidth}X}
    \toprule Método & Procedimiento \\ \midrule
    Selección hacia adelante & comienza sin atributos y agrega el que más mejora validación \\
    Eliminación hacia atrás & comienza con todos y retira el menos útil \\
    RFE & ajusta, ordena, elimina los menos importantes y repite \\
    \bottomrule
  \end{tabularx}
  \begin{block}{Ejemplo}
    Con selección hacia adelante, velocidad obtiene $F_1=0{,}71$; al agregar lluvia, $F_1=0{,}78$; agregar hora mantiene $F_1=0{,}78$. Si la regla exige una mejora mínima de 0,01, se conservan dos atributos.
  \end{block}
  \begin{alertblock}{Cuidados}
    Son costosos, no garantizan el mejor subconjunto y pueden sobreajustar la validación. La búsqueda se realiza en el ciclo interno de validación.
  \end{alertblock}
\end{frame}

% 11
\begin{frame}{Métodos embebidos}
  \footnotesize
  La selección ocurre durante el ajuste del modelo.
  \begin{tabularx}{\textwidth}{>{\bfseries}p{.2\textwidth}p{.3\textwidth}X}
    \toprule Método & Mecanismo & Ejemplo de selección \\ \midrule
    Lasso o L1 & produce algunos coeficientes iguales a cero & descarta entradas con coeficiente cero \\
    Elastic Net & combina L1 y L2 & estabiliza grupos correlacionados \\
    Árbol & utiliza atributos para construir divisiones & conserva variables presentes en ramas válidas \\
    \bottomrule
  \end{tabularx}
  \begin{block}{Ejemplo ilustrativo}
    Tras estandarizar, una regresión logística con L1 produce: velocidad $=-1{,}20$, lluvia $=0{,}70$, hora $=0$ e identificador $=0$. Bajo la regla definida, conserva velocidad y lluvia.
  \end{block}
  \begin{alertblock}{Cuidados}
    Penalización y escala se ajustan dentro de cada fold. Variables correlacionadas pueden sustituirse de forma inestable y un cero no demuestra irrelevancia causal.
  \end{alertblock}
\end{frame}

% 12
\begin{frame}{Reducción dimensional y PCA}
  \small
  La selección conserva columnas originales; la reducción construye nuevas coordenadas $Z=g(X)$ con $q<p$.
  \vspace{.35em}

  En PCA, los componentes son direcciones ortogonales ordenadas por varianza. Si los autovalores son $(4{,}2;2{,}1;0{,}7)$:
  \[
    \mathrm{PVE}_1=4{,}2/7=60\%,\qquad
    \mathrm{PVE}_{1:2}=6{,}3/7=90\%.
  \]
  \begin{alertblock}{Límite}
    Conservar 90\% de la varianza no garantiza conservar 90\% de la capacidad predictiva. Una señal rara puede encontrarse en una dirección de baja varianza.
  \end{alertblock}
\end{frame}

% 13
\begin{frame}{El pipeline completo es la unidad evaluada}
  \footnotesize
  \begin{center}
  \begin{tikzpicture}[node distance=3.5mm,scale=.64,transform shape]
    \node[datum] (d) {Datos}; \node[process,right=of d] (i) {Imputación};
    \node[process,right=of i] (c) {Codificación}; \node[process,right=of c] (e) {Escala};
    \node[process,right=of e] (s) {Selección\br o PCA}; \node[decision,right=of s] (m) {Modelo};
    \draw[flow] (d)--(i); \draw[flow] (i)--(c); \draw[flow] (c)--(e); \draw[flow] (e)--(s); \draw[flow] (s)--(m);
  \end{tikzpicture}
  \end{center}
  En cada partición:
  \begin{enumerate}
    \item ajustar transformaciones únicamente con entrenamiento;
    \item aplicar esos parámetros a validación;
    \item ajustar el modelo con la representación transformada;
    \item medir sobre los casos reservados.
  \end{enumerate}
  \begin{block}{Inferencia}
    Se reutilizan medias, categorías, columnas seleccionadas y componentes aprendidos; no se vuelven a ajustar.
  \end{block}
\end{frame}

% 14
\begin{frame}{Caso integrado: anticipar un bloqueo}
  \scriptsize
  Un vehículo debe estimar, antes de ingresar, si un tramo quedará bloqueado en los próximos 15 minutos.
  \begin{tabularx}{\textwidth}{p{.2\textwidth} p{.36\textwidth} X}
    \toprule Datos brutos & Características candidatas & Control necesario \\ \midrule
    fecha y hora & seno/coseno de hora; día hábil & calendario disponible \\
    lluvia reciente & acumulado causal; indicador intenso & corte en el instante de decisión \\
    velocidades previas & media y tendencia anteriores & excluir mediciones futuras \\
    tramo y red & tipo de vía; centralidad; desvío & política para tramos nuevos \\
    \bottomrule
  \end{tabularx}
  \begin{block}{Actividad}
    Proponer una interacción útil, una variable que causaría fuga y una prueba para detectar una diferencia entre entrenamiento e inferencia.
  \end{block}
\end{frame}

\end{document}
